\section{The actual full-Euler numerical class of the same ample model ray}
\label{sec:native-ample-model-numerical-classes}

\begin{theorem}[The original power embedding derives the true numerical scaling]
\label{thm:native-ample-model-numerical-class}
\leanok
\uses{thm:native-pulled-projective-numerical-powers,thm:native-numerical-generated-line-sums,lem:native-numerical-line-tensor-power}
Let the actual original first pulled projective twist be genuinely isomorphic
to $L^{\otimes r}$ on the same original projective scheme. Its intrinsic real
numerical class, defined by all original closed-curve full-Euler degree tests,
is exactly $r[L]$. Consequently, for $r>0$, coefficient $\theta/r$ on the same
original first-line ray is exactly $\theta[L]$, for every real $\theta$.
In particular every original power $L^{\otimes d}$ has exactly its original
numerical class at model coefficient $d/r$.
\end{theorem}
\begin{proof}
The true numerical quotient respects genuine original whole line isomorphisms.
Apply this to the actual geometric power identification, then use the proved
additivity of the same full-Euler numerical class on actual explicit tensor
powers. Scalar multiplication by $\theta/r$ and the genuine nonzero integer
$r$ give the stated equality. No desired numerical compatibility premise is
supplied, and no choice of a new numerical quotient is made.
\end{proof}

\begin{theorem}[Every original closed curve retains its actual full Euler degree]
\label{thm:native-ample-model-curve-degree}
\leanok
\uses{thm:native-ample-model-numerical-class,thm:native-pulled-projective-numerical-powers}
On every actual original closed integral curve, the true numerical degree of
this scaled model ray is $\theta$ times the actual full-Euler degree of the
original restricted line $L$. This includes every real coefficient and every
original curve test.
\end{theorem}
\begin{proof}
Apply the actual original closed-curve linear functional to the preceding
intrinsic numerical identity. Its genuine line-class formula is exactly the
full-Euler degree of the original restricted line. Linearity proves the stated
identity. This identifies the fixed ample-line ray; it does not prove
independence of arbitrary numerically equivalent line representatives or global
finite dimensionality of the numerical space.
\end{proof}
